1つ目は，日本の電力供給が現在も火力発電に大きく依存しており，燃料価格や国際情勢の影響を受けやすいからである．日本は石油，石炭，天然ガスなどの化石燃料の多くを海外からの輸入に頼っている．資源エネルギー庁によると，2023年度の日本のエネルギー自給率は15.2\%にとどまり，発電電力量においても火力発電が68.6%を占めている．そのため，国際情勢の悪化や燃料価格の上昇が，電気料金の不安定化につながりやすい構造になっている．原子力発電を一定程度活用することで，化石燃料への依存を減らし，電力供給の安定性を高めることができると考える．
2つ目は，今後さらに多くの電力が必要になると考えられるからである．近年はAIの利用拡大により，AIを動かすためのデータセンター需要が高まっている．IEAは，世界のデータセンターの電力消費量が2030年までに約945TWhへ倍増すると予測しており，これは現在の日本全体の電力消費量をやや上回る規模である．また，資源エネルギー庁も，DXやGXの進展により電力需要の大幅な増加が見込まれるとしており，データセンターや半導体工場に対して，安定的かつ国際競争力のある価格で脱炭素電源を供給する必要があると述べている．
このような将来を考えると，天候に左右される再生可能エネルギーだけで大量の電力を安定的にまかなうことは難しい場面もある．特に日本は国土が狭く，山地も多いため，発電量に対して広い土地を必要とする太陽光発電や風力発電には限界がある．その点，原子力発電は比較的小さい面積で大きな電力を安定して生み出せるため，日本のような国にとって重要な選択肢である．
もちろん，原子力発電には事故リスクや放射性廃棄物の処理といった大きな課題がある．そのため，無条件に推進するのではなく，福島第一原発事故の教訓を踏まえた安全対策，避難計画，地域住民への説明，廃棄物処理の方針を徹底することが前提である．そのうえで，安定した電力供給，電気料金の安定，脱炭素，将来のAI・データセンター需要への対応を考えると，日本は原子力発電所の建設・活用を前向きに進めるべきだと考える．